% TOP_PROBLEM: {"id": 9400081, "title": "Bombieri–Lang conjecture", "category_id": 1, "category": {"id": 1, "name": "number_theory", "display_name": "Number Theory", "description": "Properties of integers, prime numbers, Diophantine equations.", "slug": "number-theory", "order_index": 1, "created_at": "2026-07-31T15:26:25.670Z"}, "status": "partially_solved"}
% FIRST_POSTED: 2026-09-25
% ENTRY_KIND: statement_only
% !TeX program = lualatex
% Definition and sources only; this file is not an LLM solution attempt.
\documentclass[11pt]{article}
\usepackage[margin=25mm]{geometry}
\usepackage{fontspec}
\setmainfont{Latin Modern Roman}
\usepackage{amsmath,amssymb,mathtools}
\usepackage{unicode-math}
\setmathfont{Latin Modern Math}
\usepackage{xurl,hyperref}
\hypersetup{colorlinks=true,urlcolor=blue}
\setcounter{secnumdepth}{0}
\setlength{\parindent}{0pt}
\setlength{\parskip}{5pt}
\setlength{\emergencystretch}{3em}
\begin{document}
\section*{\texorpdfstring{Bombieri–Lang conjecture}{Bombieri–Lang conjecture}}\label{9400081}
\subsection{Definitions and mathematical statement}
For a smooth projective variety $X$ of dimension $d$, general type means that its canonical divisor $K_X$ is big, or equivalently its Kodaira dimension $\kappa(X)$ is $d$. For a general variety, this terminology is interpreted through a smooth projective birational model in the source's convention. The assertion is
\[
 \overline{X(K)}^{\rm Zar}\subsetneq X
\]
for every number field $K$ and every general-type $X/K$ in the chosen scope. Zariski closure is the smallest algebraic closed subset containing the rational points. Non-density is weaker than finiteness in dimensions greater than one.
\par\smallskip\textit{Formulation and concept references:} \href{https://doi.org/10.1090/S0273-0979-1986-15426-1}{[S1]}.

\subsection{Short English statement}
Must rational points on a number-field variety of general type lie in a proper algebraic subset?
\subsection{Sources}
\begin{itemize}
\item[S1]S. Lang, Bull. Amer. Math. Soc. 14(2):159-205, 1986.
\url{https://doi.org/10.1090/S0273-0979-1986-15426-1}.
\textit{Formulation source.}
\item[S2]Recent progress on the geometric Bombieri–Lang conjecture.
\url{https://arxiv.org/html/2607.02165v1}.
\textit{Further reference; not a proof endorsement.}
\item[S3]A counting argument for the geometric Bombieri-Lang conjecture on ramified covers of abelian varieties.
\url{https://arxiv.org/abs/2511.17010}.
\textit{Further reference; not a proof endorsement.}
\end{itemize}
\end{document}
